\documentclass[11pt,a4paper]{article}

\usepackage[utf8]{inputenc}
\usepackage[T1]{fontenc}
\usepackage{lmodern}
\usepackage{amsmath,amssymb,amsthm}
\usepackage{booktabs}
\usepackage[margin=2.6cm]{geometry}
\usepackage{microtype}
\usepackage{enumitem}
\usepackage[hidelinks]{hyperref}
\usepackage{xcolor}
\usepackage{listings}

\lstset{
  basicstyle=\ttfamily\small,
  breaklines=true,
  columns=fullflexible,
  keepspaces=true,
  frame=single,
  framesep=3pt,
  rulecolor=\color{gray!50},
  backgroundcolor=\color{gray!5}
}

\newtheorem{lemma}{Lemma}
\newtheorem{theorem}{Theorem}
\theoremstyle{remark}
\newtheorem*{remark}{Remark}

\newcommand{\Eone}{E_1}
\newcommand{\lam}{\lambda}
\newcommand{\R}{\mathbb{R}}
\newcommand{\floor}[1]{\left\lfloor #1 \right\rfloor}
\DeclareMathOperator{\rad}{rad}

\title{Fifty thousand certified decimal digits of\\ the Golomb--Dickman constant}

\author{Konstantin N. Pozdnyakov\\
\small Surgut, Khanty--Mansi Autonomous Okrug --- Yugra, Russia\\
\small \texttt{fix@inbox.ru}}

\date{17 November 2026}

\begin{document}

\maketitle

\begin{abstract}
We describe a rigorous interval-arithmetic certificate for the first
50\,000 decimal digits of the Golomb--Dickman constant
\[
  \lam=\int_0^\infty e^{-u-\Eone(u)}\,du
      =0.624329988543550870992936383100837244\ldots,
\]
where $\Eone$ is the exponential integral. The certificate consists of
integers $\Lambda_-<\Lambda_+$ and a scale $S=166\,293$ such that
$\lam\in[\Lambda_-/2^S,\Lambda_+/2^S]$ and
\[
  \floor{10^{50000}\Lambda_-/2^S}
  =\floor{10^{50000}\Lambda_+/2^S}.
\]
Consequently the digits are \emph{proved by truncation}, not merely
computed in floating point. The interval width is
$10^{-50042.3574}$, which leaves more than $42$ guard digits. The
computation evaluates
$\lam=\int_{\R}G(t)\,dt$ by the trapezoidal rule after the
Takahashi--Mori substitution, with explicit tail and discretisation
bounds in the spirit of Trefethen--Weideman and with every node enclosed
in an Arb ball. We give the error budget, the certificate format, a
checker that uses only exact integer arithmetic, and a b-file for OEIS
A084945 with terms $n=0,\ldots,49999$. We also report measurements on the
resulting digit string and two extensions of OEIS A225242. The
40\,000-digit certificate of the earlier run is kept unchanged as the
previous result. The trust base
is the correctness of Arb together with the cited theorems; this is not a
formalisation in a proof assistant.
\end{abstract}

\tableofcontents

\section{Introduction}

The Golomb--Dickman constant $\lam$ is the mean value of the largest
prime factor of a random integer, normalised by the integer, and the
limiting density of integers whose largest prime factor is at most
$n^{1/u}$. Dickman introduced the associated function in 1930
\cite{dickman1930}; the constant is the integral of
$F(u)=e^{-u-\Eone(u)}$ over $u\ge0$. It appears in the analysis of
random permutations, in the collision probability of hashing, and in
number-theoretic sieving; see \cite{debruijn1951} for the classical
asymptotic analysis.

Numerical values are collected in OEIS A084945 \cite{oeisA084945} and on
the MathWorld page \emph{Golomb--Dickman Constant Digits}
\cite{mathworld}. Before the work reported here the OEIS entry carried a
b-file with terms $n=0,\ldots,101$, synthesised from the entry itself,
and MathWorld quoted 15\,000 digits (2013). The present computation
extends the verified decimal expansion to 50\,000 digits and supplies a
certificate that makes the extension independent of any particular
floating-point implementation. The 40\,000-digit run of the earlier
version of this note is retained in the repository as the previous
published result; its certificate and digits are unchanged.

The contribution is deliberately narrow. We do not claim a new
algorithm: the method is the exponentially convergent trapezoidal rule
\cite{TW2014} applied after a substitution that removes the double
exponential decay, with the exponential integral evaluated by a
Stieltjes continued fraction \cite{stieltjes1894}. The contribution is
(a) a complete error budget at $D=50\,000$ digits, (b) an integer
certificate whose verification needs only exact arithmetic, and (c) a
reproducible artefact set, including checkpoints and a dependency-free
verifier.

\paragraph{Notation.}
$\Eone(u)=\int_u^\infty e^{-s}s^{-1}ds$; $\gamma$ is Euler's constant;
$D$ is the number of certified decimal digits, $G$ the number of guard
digits, and $S$ the binary scale of the certificate.

\section{The integral and the substitution}

Set
\begin{equation}
  F(u)=e^{-u-\Eone(u)},\qquad \lam=\int_0^\infty F(u)\,du .
\end{equation}
For $u>0$ one has $0<F(u)<\min(e^\gamma u,\,e^{-u})$, and $F$ extends to
an entire function of $u$ through
\begin{equation}
  F(u)=e^{\gamma}u\exp\!\Bigl(-2u+\sum_{k\ge2}\frac{(-u)^k}{k\,k!}\Bigr),
  \label{eq:Fseries}
\end{equation}
which follows from the classical expansion of $\Eone$
\cite[5.1.11]{AS1964}.

The integrand decays like $e^{-u}$ on the right and like $u$ on the
left, so the trapezoidal rule converges only algebraically. The
Takahashi--Mori substitution
\begin{equation}
  u=u(t)=\exp\bigl(t-e^{-t}\bigr),\qquad t\in\R,
  \label{eq:subst}
\end{equation}
is an increasing bijection $\R\to(0,\infty)$ with
$du=u(t)(1+e^{-t})\,dt$. Hence
\begin{equation}
  \lam=\int_{\R}G(t)\,dt,\qquad
  G(t)=F(u(t))\,u(t)\,(1+e^{-t}),
  \label{eq:G}
\end{equation}
and $G$ decays double exponentially at both ends,
$G(t)=O(e^{-e^{|t|}})$.

Two elementary tail identities are used below. For every $t_0$,
\begin{equation}
  \int_{t_0}^{\infty}u\,e^{-u}(1+e^{-t})\,dt=e^{-u(t_0)},\qquad
  \int_{-\infty}^{t_0}e^{\gamma}u^2(1+e^{-t})\,dt
  =\tfrac12 e^{\gamma}u(t_0)^2 .
  \label{eq:tails}
\end{equation}
They follow by the substitution $u=u(t)$ and the bounds of
\eqref{eq:Fseries}.

\section{Quadrature with explicit bounds}

Fix a step $h>0$ and integers $K_L,K_R$ with $u(K_Rh)\ge1$. Write
\begin{equation}
  Q_h=h\sum_{k=-K_L}^{K_R}G(kh).
\end{equation}

\begin{lemma}[Truncation]
\label{lem:tail}
With $G$ as in \eqref{eq:G},
\begin{equation}
  0\le h\!\!\sum_{k>K_R}\!G(kh)\le e^{-u(K_Rh)},\qquad
  0\le h\!\!\sum_{k<-K_L}\!G(kh)\le \tfrac12 e^{\gamma}u(-K_Lh)^2 .
\end{equation}
\end{lemma}

\begin{proof}
By \eqref{eq:Fseries}, $0<G\le g_R:=ue^{-u}(1+e^{-t})$ and
$0<G\le g_L:=e^\gamma u^2(1+e^{-t})$. The function $g_R$ decreases for
$u\ge1$, so $hg_R(kh)\le\int_{(k-1)h}^{kh}g_R$ and the right tail is at
most $\int_{K_Rh}^\infty g_R=e^{-u(K_Rh)}$ by \eqref{eq:tails}. For
$g_L$ one has $\frac{d}{dt}\log g_L=2(1+e^{-t})-e^{-t}/(1+e^{-t})>0$, so
$hg_L(kh)\le\int_{kh}^{(k+1)h}g_L$ and the left tail is at most
$\frac12e^\gamma u(-K_Lh)^2$.
\end{proof}

\begin{lemma}[Discretisation]
\label{lem:disc}
Let $d=7/5$, $c=\cos d>0$, and suppose $G$ is analytic in the strip
$|\operatorname{Im}z|<d$ with
$\int_\R|G(x+iy)|\,dx\le M$ there. Then
\begin{equation}
  \Bigl|\lam-Q_h\Bigr|
  \le B_{\mathrm{disc}}+\tau_R+\tau_L,
  \qquad
  B_{\mathrm{disc}}=\frac{2M}{e^{2\pi d/h}-1},
\end{equation}
where $\tau_R,\tau_L$ are the tail bounds of Lemma~\ref{lem:tail}.
\end{lemma}

\begin{proof}
The quadrature error is bounded by the Trefethen--Weideman trapezoidal
rule estimate \cite[Thm.~5.1]{TW2014}, and the two tails by
Lemma~\ref{lem:tail}. The decay required in the strip follows from the
estimates below.
\end{proof}

For $G$ one may take
\begin{equation}
  M=M_1+M_2+M_3,
  \label{eq:M}
\end{equation}
with
\begin{align}
  M_1&=C_1e^{-2c}\Bigl(\frac1c+\frac1{4c^2}\Bigr),
  &C_1&=e^{\gamma+2+(e-2)/2},\\
  M_2&=4e^{\gamma}e^{4}\exp\!\Bigl(2e^2+\tfrac12e^{e^2}\Bigr),\\
  M_3&=\frac{e^{\,e^{\,e^{-2}-2}}\,(1+e^{-2})\,e^{-\kappa'e^2}}{\kappa'},
  &\kappa'&=\cos(\theta_{\max})\,e^{-e^{-2}},
\end{align}
and $\theta_{\max}=d+e^{-2}\sin d<\pi/2$. The three terms bound the
strips $t\le0$, $0\le t\le2$ and $t\ge2$ respectively. The middle term
is deliberately crude and gives $M\approx10^{360}$; this is harmless,
because the discretisation error at $h\approx7.35\times10^{-5}$ is
smaller than $10^{-51594}$.

\paragraph{The constant $M_3$.}
The factor $e^{e^{e^{-2}-2}}$ is the exponential of the bound
$e^{e^{-2}-2}=1/e^{2-e^{-2}}$ for $|E_1(u)|$ in the right strip.
Let $t=x+iy$ with $x\ge2$ and $|y|\le d$, put $w=t-e^{-t}$ and
$u=e^{w}$. Then
\[
  |\operatorname{Im}w|=|y+e^{-x}\sin y|
  \le d+e^{-2}\sin d=\theta_{\max}<\pi/2,
\]
so $|\arg u|\le\theta_{\max}$ and
$\operatorname{Re}u=|u|\cos(\operatorname{Im}w)\ge|u|\cos\theta_{\max}$.
For $\operatorname{Re}z>0$ one has
$E_1(z)=e^{-z}\int_0^\infty e^{-s}/(z+s)\,ds$, hence
$|E_1(z)|\le e^{-\operatorname{Re}z}/|z|$ because
$|z+s|\ge|z|$ for $s\ge0$. On this strip
$|u|=e^{x-e^{-x}\cos y}\ge e^{2-e^{-2}}$, so
\[
  |E_1(u)|\le\frac{1}{|u|}\le e^{e^{-2}-2},
  \qquad\text{and hence}\qquad
  |e^{-E_1(u)}|\le e^{e^{e^{-2}-2}} .
\]
Together with $|1+e^{-t}|\le1+e^{-2}$ and
$|e^{-u}|=e^{-\operatorname{Re}u}\le e^{-|u|\cos\theta_{\max}}$ this
yields
\[
  |G(t)|\le e^{e^{e^{-2}-2}}(1+e^{-2})\,|u|\,
  e^{-|u|\cos\theta_{\max}} .
\]
Put $\phi(x)=x-e^{-x}\cos y$, so that $|u|=e^{\phi(x)}$ and
$\phi'(x)=1+e^{-x}\cos y\ge1$. Since $\phi(2)\ge2-e^{-2}$,
\[
  \int_2^\infty|G(x+iy)|\,dx
  \le e^{e^{e^{-2}-2}}(1+e^{-2})
  \int_{e^{2-e^{-2}}}^\infty e^{-v\cos\theta_{\max}}\,dv
  =\frac{e^{e^{e^{-2}-2}}(1+e^{-2})e^{-\kappa'e^2}}
        {\cos\theta_{\max}} .
\]
Because $\kappa'=\cos\theta_{\max}e^{-e^{-2}}\le\cos\theta_{\max}$,
the right-hand side is at most $M_3$ as defined above; the factor
$e^{e^{-2}}$ between $1/\cos\theta_{\max}$ and $1/\kappa'$ is absorbed
into the bound. Numerically $M_1=195.0281\ldots$ and
$M_3=31.8556\ldots$, while $M_2=2.4595654853\ldots\times10^{360}$
dominates, so $\log_{10}M=360.3908583900\ldots$ as used below.

\section{Enclosing the nodes}

Every value $G(kh)$ is enclosed in an Arb ball \cite{johansson2017}, and
the radius of that ball is checked to be below $2^{-E_{\mathrm{bits}}}$;
the check is part of the postcondition of the node routine, not an
heuristic. Arb provides midpoint--radius interval arithmetic with
directed rounding, so each elementary operation widens the enclosure in
a certified way.

Three independent representations of $\Eone$ are used, and the routine
chooses the cheapest one that is guaranteed to reach the target radius.

\begin{enumerate}[leftmargin=1.6em]
\item \emph{Power series.} From \eqref{eq:Fseries} with
  $a_j=(-u)^j/(j\,j!)$, the ratio test gives
  $|a_{j+1}/a_j|=uj/(j+1)^2\le u/(j+1)$. If $K+1\ge2u$, the tail is at
  most $2|a_{K+1}|$; the finite part is summed in Arb and the tail is
  added as a symmetric interval.
\item \emph{Stieltjes continued fraction.} For $z>0$,
  \begin{equation}
    e^z\Eone(z)=\frac1z\,g(1/z),\qquad
    g(x)=\int_0^\infty\frac{e^{-s}}{1+sx}\,ds .
  \end{equation}
  The measure $e^{-s}ds$ has moments $m_j=j!$, and Carleman's condition
  holds, so the Stieltjes transform is determined by its moments. The
  J-fraction obtained from the Laguerre recurrence compresses to the
  S-fraction with coefficients $\alpha_{2j-1}=\alpha_{2j}=j>0$. For a
  positive Stieltjes fraction, consecutive convergents bracket the
  value; if $D_n$ bounds $|f_n-f_{n-1}|$ in interval arithmetic, then
  $f_n+[-D_n,D_n]$ contains $e^z\Eone(z)/z$ and hence the node ball.
\item \emph{Asymptotic expansion.} From
  $e^u\Eone(u)=\int_0^\infty e^{-s}/(u+s)\,ds$ and the finite geometric
  identity,
  \begin{equation}
    e^u\Eone(u)=\sum_{j=0}^{K-1}\frac{(-1)^j j!}{u^{j+1}}+R_K,
    \qquad |R_K|\le\frac{K!}{u^{K+1}} .
  \end{equation}
  This is used only where $u$ is large enough for the bound to be below
  the target radius.
\end{enumerate}

A fourth branch handles nodes whose value is provably negligible: from
$0<F(u)<\min(e^\gamma u,e^{-u})$ one gets
$0<G(t)\le e^\gamma u^2(1+e^{-t})$ for $t<0$ and
$0<G(t)\le ue^{-u}(1+e^{-t})$ for $t\ge0$. The branch returns $[0,U]$
only after Arb has certified $U<\varepsilon$; otherwise the regular
branch is used.

\begin{remark}
The choice of branch, the number of series terms and the number of
continued-fraction levels are made using floating-point estimates.
Those estimates only select resources. The enclosure and the radius
check are always performed in Arb; if the requested accuracy is not
reached, the precision is raised and the node is recomputed, or the
program aborts. A node is never accepted on the strength of a
floating-point estimate.
\end{remark}

\section{The certificate and its error budget}

Let $S=E_{\mathrm{bits}}+64$. Each accepted node contributes an interval
$[g_k^-,g_k^+]$ with $G(kh)\in[g_k^-,g_k^+]$ and
$\rad<2^{-E_{\mathrm{bits}}}$, expressed in integer units of $2^{-S}$.
The computation stores, for every chunk of nodes, the exact integer sums
of the lower and upper endpoints. Summing the chunks gives integers
$L_{\mathrm{sum}}$ and $H_{\mathrm{sum}}$ with
\begin{equation}
  h\,2^{-S}L_{\mathrm{sum}}\le Q_h\le h\,2^{-S}H_{\mathrm{sum}} .
\end{equation}
The certificate then records
\begin{equation}
  \Lambda_-=\floor{hL_{\mathrm{sum}}}-B,\qquad
  \Lambda_+=\lceil hH_{\mathrm{sum}}\rceil+C,\qquad
  C=T_R+T_L+B,
  \label{eq:cert}
\end{equation}
where
\begin{equation}
  B=\floor{B_{\mathrm{disc}}\,2^S}+1,\qquad
  T_R=\floor{\tau_R\,2^S}+1,\qquad
  T_L=\floor{\tau_L\,2^S}+1 .
\end{equation}
Thus every quantity that enters $\Lambda_\pm$ is either an exact integer
or an outward-rounded bound; no fitted constant is involved.

Table~\ref{tab:params} lists the parameters of the published run and the
resulting bounds. The numbers in the last block are reproduced
independently by \texttt{verification/level\_b.py}.

\begin{table}[ht]
\centering
\small
\begin{tabular}{@{}llr@{}}
\toprule
Quantity & Symbol & Value \\
\midrule
certified decimal digits        & $D$              & 50\,000 \\
guard digits                    & $G$              & 30 \\
node radius exponent            & $E_{\mathrm{bits}}$ & 166\,229 \\
certificate scale               & $S$              & 166\,293 \\
step                            & $h$              & $80\,846\,368/2^{40}$ \\
left/right node counts          & $K_L,K_R$        & 149\,075,\ 158\,504 \\
chunk size                      & ---              & 100 \\
chunks / nodes                  & ---              & 3\,076 / 307\,580 \\
\midrule
strip half-width                & $d$              & $7/5$ \\
bound $\log_{10}M$              & ---              & $360.3909$ \\
discretisation bound            & $\log_{10}B_{\mathrm{disc}}$ & $-51594.8080$ \\
right tail bound                & $\log_{10}\tau_R$ & $-50043.8007$ \\
left tail bound                 & $\log_{10}\tau_L$ & $-50045.7499$ \\
\midrule
lower correction                & $B$              & 1 \\
right tail correction           & $T_R$            & 2\,400\,727\,111\,627\,803 \\
left tail correction            & $T_L$            & 26\,985\,556\,367\,958 \\
upper correction                & $C=T_R+T_L+B$    & 2\,427\,712\,667\,995\,762 \\
\midrule
interval width                  & $\log_{10}(\Lambda_+-\Lambda_-)-S\log_{10}2$ & $-50042.357438746905$ \\
guard digits beyond $10^{-D}$   & ---              & $42.3574$ \\
\bottomrule
\end{tabular}
\caption{Parameters and error budget of the 50\,000-digit run. All
integer quantities are exact.}
\label{tab:params}
\end{table}

The discretisation bound is about $1559$ orders of magnitude below what
is needed, i.e.\ below $10^{-(D+G+6)}$, and each tail is about $6$ to
$8$ orders below its target $10^{-(D+G+8)}$. The
width is dominated by the accumulated widths of the saved node
enclosures: the chunk-sum width is $10^{-50042.3737}$, while the
outward corrections $B$ and $C$ add $2^{51.1}$ units of $2^{-S}$ to a
total of $2^{55.89}$. The analytic error terms do not control the
guard-digit count; the interval arithmetic does.

\section{Verification}

The certificate is a three-line text file: a comment line naming the
constant and the method, a comment line with $S$ and $\Lambda_-$, a
comment line with $\Lambda_+$, and a line containing
$0.d_1d_2\cdots d_D$. Verification needs only exact integer arithmetic.

\begin{enumerate}[leftmargin=1.6em]
\item Parse $S,\Lambda_-,\Lambda_+$ and the digit string.
\item Check $\floor{10^D\Lambda_-/2^S}=\floor{10^D\Lambda_+/2^S}$;
  this is the statement that the interval does not cross a decimal cell.
\item Check that the common value, written with $D$ leading zeros if
  necessary, equals the digit string.
\item Check that the standalone digits file equals the certificate line.
\item Check that the b-file has exactly $D$ terms, indexed
  $0,\ldots,D-1$, each equal to the corresponding digit, in plain ASCII
  with LF line endings.
\item Check that the first 102 digits equal the published prefix of OEIS
  A084945.
\end{enumerate}

Steps 1--6 are implemented in \texttt{40000/verify\_40000.py} with the
Python standard library only, and run in about $0.1$\,s. The repository
also contains a secondary audit that re-reads all $1\,209$ checkpoints,
recomputes the exact integer sums, re-derives $B$, $T_R$, $T_L$ from the
formulas of Section~3, and checks that the published endpoints are
exactly reproduced
(\texttt{checks/verify\_checkpoints.py}). A third script recomputes
selected nodes from scratch in \texttt{decimal}
(\texttt{checks/verify\_nodes\_decimal.py}); a fourth re-implements the
three representations of $\Eone$ and recomputes all nodes in Arb
(\texttt{checks/verify\_checkpoints\_arb.py}).

The last script is the only check that re-does the Arb evaluation
independently of the producing program. It was not executed on the
machine used to assemble this package, because an Application Control
policy blocked the Arb shared libraries; it is included for a host
without that restriction. Consequently the certificate presented here is
conditional on the producing program having computed every ball
correctly, on Arb being correct, and on the elementary lemmas of
Section~3 and Section~4. It is not a formally verified proof.

The repository organises the checks in three levels, so that a reader can
pay exactly as much as the question they care about costs
(\texttt{verification/}). Level~A verifies the certificate and the digit
strings with exact integer arithmetic and no external dependency
(\texttt{verify\_40000.py}). Level~B audits the parameters, the
discretisation bound and the two tails, and recomputes a sample of chunks
from the formulas (\texttt{level\_b.py}); saved checkpoints are treated as a
cache, never as evidence. Level~C recomputes the leading digits with a
different representation of the integral, a different quadrature rule and a
different library (\texttt{independent\_quadrature.py}, \texttt{mpmath}).
Level~C is floating point, hence numerical evidence rather than a proof; it
reproduces the first $1\,000$ digits in our runs.

\section{The digit string}

Table~\ref{tab:digits} collects measurements on the first $40\,000$
digits of the certified $50\,000$-digit string. They are reproducible
with \texttt{40000/checks/verify\_sequences.py} in about $12$\,s. None of
these observations is used in the proof; they document the object and
provide independent consistency checks against published statements.

\begin{table}[ht]
\centering
\small
\begin{tabular}{@{}p{0.42\textwidth}p{0.5\textwidth}@{}}
\toprule
Quantity & Value \\
\midrule
digit frequencies $0,\ldots,9$ &
  $3993,3950,3969,4076,3993,3972,3989,4077,3986,3995$ \\
$\chi^2$ (9 degrees of freedom) & $4.10$ \\
longest run of one digit & 5 (\texttt{00000}, \texttt{11111}, \texttt{22222}, \texttt{44444}) \\
longest repeated block & \texttt{1891192031} at positions 28\,793 and 34\,393 \\
longest square $XX$ & \texttt{54395439} at position 30\,398 \\
longest palindrome & \texttt{05224342250} at position 11\,751 \\
prime substrings of length $\le32$ & 66\,816 (deterministic to 24 digits) \\
prime truncations up to 3000 digits & lengths 6, 27, 57, 60, 1659, 2508 \\
\bottomrule
\end{tabular}
\caption{Measurements on the first $40\,000$ digits of the certified
$50\,000$-digit string. The
$\chi^2$ value is consistent with uniformity but proves nothing about
normality, which remains open.}
\label{tab:digits}
\end{table}

Two statements in OEIS A225242 \cite{oeisA225242} are extended. Writing
$a(n)$ for the first position at which the digit $n$ appears $n$ times in
a row, the published entry recorded $a(5)>15\,000$; the strings
\texttt{55555} and \texttt{666666} do not occur among the $50\,000$
digits, so $a(5),a(6)>50\,000$. Neither \texttt{0123456789} nor
\texttt{9876543210} occurs either, and the longest run of a single digit
has length $5$. This is a conditional statement: it is
valid for the certified string, hence for $\lam$ if the certificate is
correct.

\section{Reproducibility and data}

The producing pipeline is \texttt{50000/golomb\_dickman\_26.py}. It requires
Python 3.11+ and \texttt{python-flint} \cite{pythonflint} (Arb); the run
takes hours of CPU and is resumable through checkpoints. The repository
ships the $3\,076$ published checkpoints, so the certificate can be
re-derived without recomputing the nodes:

\begin{lstlisting}
python -m pip install python-flint
cd 50000
python ../verification/verify_certificate.py \
    --cert golomb_dickman_50000_certified.txt \
    --sha256-manifest SHA256SUMS.txt
python ../verification/run_all.py        # levels A and B, about 1 min
python golomb_dickman_26.py certify --digits 50000   # hours
\end{lstlisting}

The command-line interface accepts
\texttt{certify --digits 50000}, \texttt{crosscheck},
\texttt{boundary-test}, \texttt{test} and \texttt{smoke}.

\paragraph{Data.}
The certificate \texttt{golomb\_dickman\_50000\_certified.txt}, the
digit file \texttt{golomb\_dickman\_50000\_digits.txt}, the OEIS b-file
\texttt{b084945.txt}, the code, the verifiers and the checkpoints are
archived in the GitHub repository
\url{https://github.com/fix27/golomb-dickman-constant/tree/main/50000}.
The 40\,000-digit certificate and b-file of the previous run are kept in
\texttt{40000/} byte-for-byte. SHA-256 checksums of the canonical files
are given in Appendix~\ref{app:sums}. An archived release with a DOI is
planned; see \texttt{40000/PUBLICATION\_CHECKLIST.md}.

\paragraph{Trust base.}
The result is conditional on:
(i) the correctness of Arb's midpoint--radius arithmetic
\cite{johansson2017};
(ii) the trapezoidal rule theorem of Trefethen and Weideman
\cite[Thm.~5.1]{TW2014};
(iii) the classical Stieltjes theorem on two-sided bounds by consecutive
convergents \cite{stieltjes1894};
(iv) the elementary lemmas of Sections~2--4, stated with proofs in the
docstring of \texttt{golomb\_dickman.py};
(v) exact integer arithmetic in Python.
No step depends on the normality of $\lam$, on unproved conjectures, or
on the numerical value of $\gamma$ beyond its certified enclosure.

\section{Conclusion}

The first 40\,000 decimal digits of the Golomb--Dickman constant are
certified by an interval of width $10^{-40042.26}$ whose endpoints give
the same $40\,000$-digit truncation. The certificate is small, the
verifier is dependency-free, and the b-file for OEIS A084945 extends the
published data from 101 to 40\,000 terms. The main open task is the
independent Arb re-audit of all nodes on an unrestricted host; the script
is part of the package.

\appendix

\section{First digits}

The certified value begins
\begin{lstlisting}
0.6243299885435508709929363831008372441796426201805292869735519024956380888551132544624602761955398688...
\end{lstlisting}
The digits above are the first 100 digits after the point; the full
string has 50\,000 digits and is given in
\texttt{50000/golomb\_dickman\_50000\_digits.txt}.

\section{Checksums}
\label{app:sums}

SHA-256 of the canonical files of the 50\,000-digit run (LF line
endings, ASCII/UTF-8 without BOM):
\begin{lstlisting}
76f4155981049003351409769505c3cec39e26dfc30624c46bd199e70ed6477a  golomb_dickman_50000_certified.txt
8f017d45ee5108bb6b976882b079279f44f8a767ce62b44bd7c54e9167ad09a0  golomb_dickman_50000_digits.txt
1e3d23e9ac3ce51990d83decd395d6bde96027f43811b94cebd5a6a4b7c19c43  b084945.txt
e2eb7d4ccca761c5c5e1f2b37075f9baea245dac86cbe962fd9dc92aa574ef71  golomb_dickman_26.py
78783a09bae599298a07415056b0ca521717c9f7ecaa1fd15cb19b4d5ab6b0e0  gd_rigorous_ckpt_50000/meta_v12.json
\end{lstlisting}
The 40\,000-digit run of the previous version, kept unchanged:
\begin{lstlisting}
8892c0715f1afbb271c9f07cc9658cea0f3ca83c2dda3ec7d23b447dc908de9c  golomb_dickman_40000_certified.txt
fbae757900983bd4fb6421c5f27bf9602af1a1bfc013e4c501bb9da2b5a65327  golomb_dickman_40000_digits.txt
c73af643198da299d2e624e3c2f69d05439efcddbbefa7cec176ac63f993bed3  b084945.txt
\end{lstlisting}

\begin{thebibliography}{99}

\bibitem{dickman1930}
K.~Dickman,
\emph{On the frequency of numbers containing prime factors of a certain
relative magnitude},
Arkiv f\"or Matematik, Astronomi och Fysik \textbf{22A} (1930), no.~10,
1--14.

\bibitem{debruijn1951}
N.~G. de~Bruijn,
\emph{The asymptotic behaviour of a function occurring in the theory of
primes},
Journal of the Indian Mathematical Society \textbf{15} (1951), 25--32.

\bibitem{AS1964}
M.~Abramowitz and I.~A. Stegun (eds.),
\emph{Handbook of Mathematical Functions with Formulas, Graphs, and
Mathematical Tables},
National Bureau of Standards, 1964.

\bibitem{TW2014}
L.~N. Trefethen and J.~A.~C. Weideman,
\emph{The exponentially convergent trapezoidal rule},
SIAM Review \textbf{56} (2014), no.~3, 385--458.

\bibitem{johansson2017}
F.~Johansson,
\emph{Arb: efficient arbitrary-precision midpoint--radius interval
arithmetic},
IEEE Transactions on Computers \textbf{66} (2017), no.~8, 1281--1292.

\bibitem{stieltjes1894}
T.-J. Stieltjes,
\emph{Recherches sur les fractions continues},
Annales de la Facult\'e des Sciences de Toulouse \textbf{8} (1894),
J1--J122.

\bibitem{oeisA084945}
OEIS Foundation Inc.,
\emph{Sequence A084945: Decimal expansion of the Golomb--Dickman
constant},
\url{https://oeis.org/A084945}.

\bibitem{oeisA225242}
OEIS Foundation Inc.,
\emph{Sequence A225242},
\url{https://oeis.org/A225242}.

\bibitem{mathworld}
E.~W. Weisstein,
\emph{Golomb--Dickman Constant Digits},
MathWorld, Wolfram Research,
\url{https://mathworld.wolfram.com/Golomb-DickmanConstantDigits.html}.

\bibitem{pythonflint}
F.~Johansson et al.,
\emph{python-flint: Python bindings for FLINT and Arb},
\url{https://github.com/flintlib/python-flint}.

\end{thebibliography}

\end{document}
